% =====================================================================================
% Figure -- the ArUco space-pin markers in situ across the study site.
% Belongs in the alignment subsection, alongside the description of the World Locking
% Tools pin network (METHODOLOGY_NOTES.md section 5).
%
% Requires in the preamble (acmart does not load it by default):
%     \usepackage{subcaption}
%
% Expected image files, relative to the main .tex:
%     figures/marker-curb.jpg      (curb face beside stair handrail)
%     figures/marker-doorpad.jpg   (service entrance pad)
%     figures/marker-pier.jpg      (base of a stone-clad pier)
%     figures/marker-corner.jpg    (brick building corner at a planting bed)
%     figures/marker-stair.jpg     (foot of an exterior stairwell)
%     figures/marker-joint.jpg     (slab joint beside a ramp wall)
%
% ASPECT RATIO -- all six sources are square (2048 x 2048), so they tile at equal height
% with no cropping. Do not substitute a non-square image into this grid without cropping
% it square first, or that row will render at a different height:
%     magick figures/NEW.jpg -gravity center -crop 2048x2048+0+0 +repage figures/NEW.jpg
%
% The marker is small in frame in every shot, which is the point -- it shows the real
% placement scale. If it reproduces too faintly in print, crop each source tighter around
% the marker rather than enlarging the panels.
%
% CAPTION CAVEAT -- the numbers below come from METHODOLOGY_NOTES.md (section 5, lines
% 122--124 and 209--216): Dictionary_5x5_250, 0.15 m space-pin markers, 0.1016 m map
% marker, 19 ArucoPinDriver instances (18 space pins + 1 map marker). Confirm these
% against the shipped scene before submitting, and confirm the six panels here are space
% pins and not the map marker.
%
% Reference in text as Figure~\ref{fig:markers}.
% =====================================================================================

\begin{figure*}[t]
  \centering
  \begin{subfigure}[b]{0.32\textwidth}
    \includegraphics[width=\linewidth]{figures/marker-curb}
    \caption{Curb face.}
    \label{fig:marker-curb}
  \end{subfigure}
  \hfill
  \begin{subfigure}[b]{0.32\textwidth}
    \includegraphics[width=\linewidth]{figures/marker-doorpad}
    \caption{Entrance pad.}
    \label{fig:marker-doorpad}
  \end{subfigure}
  \hfill
  \begin{subfigure}[b]{0.32\textwidth}
    \includegraphics[width=\linewidth]{figures/marker-pier}
    \caption{Pier base.}
    \label{fig:marker-pier}
  \end{subfigure}

  \vspace{2pt}

  \begin{subfigure}[b]{0.32\textwidth}
    \includegraphics[width=\linewidth]{figures/marker-corner}
    \caption{Building corner.}
    \label{fig:marker-corner}
  \end{subfigure}
  \hfill
  \begin{subfigure}[b]{0.32\textwidth}
    \includegraphics[width=\linewidth]{figures/marker-stair}
    \caption{Stairwell foot.}
    \label{fig:marker-stair}
  \end{subfigure}
  \hfill
  \begin{subfigure}[b]{0.32\textwidth}
    \includegraphics[width=\linewidth]{figures/marker-joint}
    \caption{Slab joint.}
    \label{fig:marker-joint}
  \end{subfigure}
  \caption{Six of the ArUco markers that register the digital twin to the physical site,
    photographed in place. Markers are 0.15\,m printed squares from the
    \texttt{5x5\_250} dictionary, with the physical side length declared explicitly to
    the tracker rather than estimated. Each marker is fixed at a surveyed position and
    drives a World Locking Tools space pin: rather than the model snapping to a single
    marker, every marker pins a known model point to a known world point, and the
    alignment is the fit over the pin network. Markers are scanned once before trials
    begin and the twin is then held in place by the locked coordinate frame, so no marker
    needs to remain in view during a trial. Placement follows a consistent rule across
    the site: markers sit flat on paving or on the vertical face of a curb, at durable
    concrete edges, corners, joints, and thresholds, low in the visual field and out of
    the line of foot traffic. The features they attach to are precisely the ones that
    reconstruct cleanly in the site scan, so the surveyed model point is unambiguous;
    keeping them off doors, signage, and vegetation also keeps them from moving between
    sessions.}
  \Description{Six square photographs of small printed black-and-white ArUco markers
    fixed to concrete surfaces around a campus site, each marker roughly the size of a
    hand and occupying a small part of the frame. In the first, a marker is stuck to the
    vertical face of a low concrete curb beside a lawn, next to the foot of a stainless
    steel stair handrail. In the second, a marker sits on the edge of a raised concrete
    pad in front of grey service doors in a brick wall, with oil stains on the pavement
    beside it. In the third, a marker lies flat on the pavement in a slab joint at the
    base of a stone-clad pier. In the fourth, a marker lies on the pavement at the
    outside corner of a red brick building where the paving meets a planting bed of
    grasses and shrubs, with a wooden board leaning against the wall. In the fifth, a
    marker sits on the concrete floor at the foot of an exterior stairwell between two
    concrete walls. In the sixth, a marker lies on the pavement alongside the sealed
    joint at the base of a sloping concrete ramp wall. In every case the marker is on
    bare, hard, permanent construction rather than on a door, a sign, or planting.}
  \label{fig:markers}
\end{figure*}
